\chapter{Related Work}
\label{chap:related-work}

This chapter situates DRIFT among four bodies of work: production memory systems for language model agents, research on the risks of memory consolidation, write-side control techniques for what enters memory in the first place, and benchmarks that treat staleness as a first-class evaluation target rather than an afterthought. Each section closes by stating precisely what the cited work does not do, which is the gap this thesis addresses.

\section{Agent Memory Systems}
\label{sec:related-memory-systems}

The last several years have produced a series of memory architectures aimed at extending a language model agent's effective context beyond what fits in a single prompt.

MemGPT introduced an operating-system-inspired design in which the agent's context window is treated as physical memory and an external store is treated as disk, with the model itself issuing function calls to page information in and out \citep{packer2023memgpt}. This design is primarily concerned with the mechanics of what enters and leaves the active context at a given moment, and it does not distinguish between information that is merely old and information that has been superseded by a later observation.

Zep constructs a temporal knowledge graph specifically for agent memory and is the closest architectural relative of the substrate used in this thesis \citep{rasmussen2025zep}. Zep records the interval during which a fact was valid, which is the same primitive DRIFT's substrate relies on. What Zep's published work does not do is compare alternative policies for acting on that temporal information: it demonstrates that recording validity intervals is useful for retrieval quality, but it does not ask which of several possible rules for using those intervals to forget, invalidate, or retain facts performs best, nor does it report a staleness metric distinct from answer accuracy. This is precisely the comparison DRIFT is built to run, using a similar temporal primitive as its starting point.

Mem0 targets production deployment with a scalable long-term memory store, emphasizing operational concerns such as latency and storage cost at scale \citep{chhikara2025mem0}. A-MEM organizes memory agentically, delegating decisions about memory structure and linking to a language model rather than fixing them in advance \citep{xu2025amem}.

All four systems are systems contributions: each proposes an architecture and demonstrates that the architecture improves some measure of retrieval quality or operational efficiency relative to a simpler baseline. None of the four isolates the forgetting rule as a variable that can be swapped while everything else about the system is held constant, and none reports an accuracy metric alongside a metric that specifically measures whether the agent's answer reflects a fact that has already been superseded. Comparing them to one another, or to any single one of them, cannot answer the question this thesis asks, because in each comparison the forgetting rule changes together with the substrate, the retrieval function, and often the underlying model.

\section{Consolidation and Its Risks}
\label{sec:related-consolidation}

A separate line of work treats memory consolidation, the periodic act of rewriting, summarizing, or pruning stored memories, as a technique that can degrade an agent's behavior rather than only improve it. Zhang et al. show that memories consolidated through repeated language-model rewriting degrade in quality over time, to the point where an agent's task utility can fall below what the same agent would achieve with no memory at all \citep{zhang2026faulty}. Their headline finding, that a capable model fails on the majority of a set of previously solved reasoning problems once its memory of the solutions has been consolidated by repeated rewriting, is a direct warning against treating consolidation as a strictly beneficial operation. This finding motivates one of the design decisions of this thesis: rather than assuming that language-model consolidation is an improvement, DRIFT treats it as a policy under test, on equal footing with every other candidate. It appears in the taxonomy of Chapter~\ref{chap:method} as policy P4, evaluated with the same metrics and the same statistical rigor applied to every other policy, rather than assumed to be a baseline improvement.

\section{Write-Side Control}
\label{sec:related-write-side}

A third line of work intervenes earlier in the pipeline, at the point where a candidate fact is either admitted to memory or rejected before it is ever stored. Wang et al. propose SAGE, which uses a von Mises-Fisher density estimator over memory embeddings to classify an incoming fact as clearly novel, clearly redundant, or uncertain, routing clearly novel facts directly into memory, discarding clearly redundant ones, and reserving a more expensive language-model merge step only for the uncertain middle case \citep{wang2026sage}. The reported benefit is substantial: a reduction of roughly three and a half times in the cost of the write-time API calls needed to maintain memory. This technique is the direct model for the novelty-gating policy evaluated in this thesis as P3. Chapter~\ref{chap:results} reports a specific and, as far as the author is aware, previously undemonstrated failure mode that follows from the design of any novelty gate, not from a particular implementation of one: a correction to an existing fact is, by construction, textually close to the fact it corrects, which places corrections in exactly the region of embedding space a novelty gate is designed to reject.

\section{Staleness as a First-Class Problem}
\label{sec:related-staleness}

A more recent line of work asks directly whether an agent can recognize that something it believes is no longer true. Chao et al. introduce STALE, a benchmark constructed to test whether an agent recognizes that a stored belief has become outdated, and identify implicit conflict, in which a new observation contradicts an old belief without explicitly announcing that it does so, as an underexplored failure mode \citep{chao2026stale}. Patel frames a closely related gap as a memory-update problem and studies remedies applied at training time rather than at inference time \citep{patel2026supersede}. Du surveys the broader field of memory mechanisms for autonomous agents and identifies the evaluation of memory revision, meaning an agent's ability to correctly update what it believes in response to new evidence, as an open frontier for the field \citep{du2026memory}.

The contribution of this thesis is complementary to, and orthogonal to, this line of work. Rather than asking whether a model can detect that its own memory is stale, or training a model to do so, this thesis asks a different and logically prior question: which forgetting policy, applied to the memory substrate before the model ever answers, prevents staleness from arising as a retrievable option in the first place, while holding the underlying model fixed. If the policy governing what memory retains can prevent a stale fact from ever being surfaced to the model, the model's own ability to detect staleness becomes a secondary line of defense rather than the primary one.

\section{Benchmarks}
\label{sec:related-benchmarks}

LoCoMo and LongMemEval are the two benchmarks most commonly used to evaluate agent memory over long conversational histories, covering multi-hop and temporal reasoning questions grounded in extended dialogue \citep{maharana2024locomo,wu2025longmemeval}. Neither benchmark annotates the value a stale memory would produce for any question in its question set, a gap examined in detail in Section~\ref{sec:benchmark-selection}. MemoryAgentBench contributes a split called FactConsolidation, constructed so that later facts in a session overwrite earlier facts sharing the same subject and predicate, which makes it directly usable for measuring whether an agent's answer reflects the current or the superseded value \citep{hu2026memoryagentbench}. MQuAKE provides knowledge-editing instances that ship both the pre-edit and the post-edit answer for each edited fact, derived from real changes to Wikidata over time rather than from synthetic construction \citep{zhong2023mquake}. GateMem studies memory governance in settings with multiple agents sharing one memory store, which is the setting most relevant to the fourth research question of this thesis, concerning whether policy rankings established under single-agent use survive multi-agent use \citep{ren2026gatemem}.

\section{Summary}
\label{sec:related-summary}

Across these four bodies of work, a consistent pattern emerges. Systems papers propose architectures and demonstrate improvements in retrieval quality or operational cost without isolating the forgetting rule as a variable. Consolidation research demonstrates that a specific technique, language-model rewriting, can harm rather than help, but does not situate that technique within a broader controlled comparison against alternatives. Write-side control research proposes an efficiency technique without characterizing its failure modes on workloads where facts are corrected rather than simply accumulated. Staleness-focused benchmark research asks whether a model can detect that its beliefs are outdated, a question that is complementary to, but distinct from, the question of which policy prevents staleness from being retrievable in the first place. The gap common to all four is the absence of a controlled comparison in which the forgetting policy is varied while every other component of the pipeline, including the substrate, the retrieval function, the extraction step, the answer model, and the judge, is held fixed. Supplying that comparison is the purpose of the framework described in the next chapter.
